% BHU PYQ -- Manually transcribed & error-corrected
\documentclass[11pt,a4paper]{article}

\usepackage[a4paper,top=2.5cm,bottom=2.5cm,left=2.8cm,right=2.8cm,headheight=14pt]{geometry}
\usepackage{amsmath,amssymb}
\usepackage[T1]{fontenc}
\usepackage[utf8]{inputenc}
\usepackage{lmodern}
\usepackage{microtype}
\usepackage{fancyhdr}
\usepackage{enumitem}
\usepackage{hyperref}
\usepackage{lastpage}
\usepackage{parskip}

\hypersetup{colorlinks=true,urlcolor=black,linkcolor=black,
  pdftitle={BPT-503: Quantum Mechanics -- BHU PYQ},pdfauthor={Banaras Hindu University}}

\pagestyle{fancy}\fancyhf{}
\renewcommand{\headrulewidth}{0.4pt}\renewcommand{\footrulewidth}{0.4pt}
\fancyhead[L]{\small   \textit{Banaras Hindu University}}
\fancyhead[R]{\small   \textit{BPT-503: Quantum Mechanics}}
\fancyfoot[C]{\small Page  \thepage\ of \pageref{LastPage}}


\newlist{parts}{enumerate}{2}
\setlist[parts,1]{label=(\alph*),leftmargin=*,itemsep=5pt,topsep=3pt}
\setlist[parts,2]{label=(\roman*),leftmargin=*,itemsep=3pt,topsep=2pt}

\newcommand{\pts}[1]{\hfill   \textit{[#1]}}


\begin{document}


\begin{center}
  {\large\bfseries BANARAS HINDU UNIVERSITY}\\[2pt]
  {
\normalsize B.Sc. (Hons.) Semester V Examination, 2013-14}\\[6pt]
  \rule{0.8\linewidth}{0.5pt}\\[6pt]
  {\large\bfseries Paper No. BPT-503: Quantum Mechanics}\\[4pt]
  {
\normalsize Time: 3 Hours\hfill Maximum Marks: 70}
\end{center}
\rule{\linewidth}{0.4pt}
\smallskip



\noindent   \textit{Instructions: Answer all questions. Symbols have their usual meanings.}
\bigskip




\noindent   \textbf{Question 1.} Answer any   \textbf{four} of the following: \pts{4 $ \times$ 5 = 20}

\begin{parts}
  \item Find the de Broglie wavelength associated with:
    
\begin{parts}
      \item an electron accelerated through a potential difference of $10\, \text{MV}$.
      \item a proton having a kinetic energy of $10\, \text{MeV}$.
    \end{parts}
    (Take rest mass energy of electron = $0.51\, \text{MeV}$ and rest mass energy of proton = $938\, \text{MeV}$).
  \item Using the uncertainty relation, explain why electrons cannot exist inside the nucleus.
  \item What are stationary states? Explain with reference to the time-dependent and time-independent Schr\"{o}dinger equations.
  \item Prove the statement: ``commuting observables have simultaneous eigenstates''.
  \item Evaluate the following commutation relations:
    
\begin{parts}
      \item $[L_x, y]$
      \item $[L_z, p_x]$
    \end{parts}
  \item Find the most probable value of the radial distance $r$ of the electron in the ground state of the Hydrogen atom.
\end{parts}

\medskip\rule{\linewidth}{0.2pt}




\noindent   \textbf{Question 2.}

\begin{parts}
  \item Describe the difference between phase velocity and group velocity for a wave packet. \pts{4}
  \item A particle is constrained to move in a potential of the form:
    \[ V(x) = 
\begin{cases} 0, & 0 < x < a \\ \infty, &  \text{otherwise} \end{cases} \]
    Obtain the values of energy $E$ and corresponding normalized eigenfunctions. Show that the eigenfunctions are mutually orthogonal. Is there any acceptable physical solution with $E=0$ or $E<0$? Justify your answer. \pts{10}
\end{parts}

\medskip\rule{\linewidth}{0.2pt}




\noindent   \textbf{Question 3.}

\begin{parts}
  \item Describe the Davisson-Germer experiment to demonstrate the wave nature of the electron. \pts{4}
  \item If a particle having energy $E > 0$ is incident from the left upon a one-dimensional square potential well defined as:
    \[ V(x) = 
\begin{cases} -V_0, & -a < x < a \\ 0, &  \text{otherwise} \end{cases} \]
    find the transmission coefficient $T$ and plot its variation qualitatively as a function of the ratio $E/V_0$. \pts{10}
\end{parts}
\hfill   \textbf{OR}

\begin{parts}
  \item Using the Schwarz inequality, give the mathematical proof of the uncertainty relation for canonically conjugate pairs. \pts{6}
  \item Find the wave function $\psi(x)$, which is the Fourier transform of $\phi(k)$, for a wave packet whose amplitude function $\phi(k)$ varies as:
    \[ \phi(k) = 
\begin{cases} g_0, & -K_0 < k < K_0 \\ 0, &  \text{otherwise} \end{cases} \]
    Calculate the probability of finding the particle in the region $-\frac{\pi}{K_0} < x < \frac{\pi}{K_0}$. \pts{8}
\end{parts}

\medskip\rule{\linewidth}{0.2pt}




\noindent   \textbf{Question 4.}

\begin{parts}
  \item What is the difficulty in normalizing a free particle wave function? How is the concept of a wave packet used to resolve this difficulty? \pts{5}
  \item Find the uncertainty in position and momentum for the Gaussian wave packet $\psi(x) = A e^{-a x^2}$. How does it obey the Heisenberg uncertainty relation? \pts{9}
\end{parts}
\hfill   \textbf{OR}

\begin{parts}
  \item The solution of the Schr\"{o}dinger equation for a free particle in one dimension is of the form:
    \[ \psi(x) = A \sin(kx) + B \cos(kx) \]
    where $A$ and $B$ are complex constants.
    
\begin{parts}
      \item Calculate the probability current density for this wave function. \pts{4}
      \item Find the relation between the coefficients $A$ and $B$, if $\psi(x)$ is an eigenfunction of the momentum operator. \pts{4}
    \end{parts}
  \item Show that the Harmonic Oscillator Hamiltonian can be written as:
    \[ H = \hbar\omega \left( a^{\dagger}a + \frac{1}{2} \right) \]
    in terms of the lowering and raising operators:
    \[ a = \sqrt{\frac{m\omega}{2\hbar}} \left( x + \frac{ip}{m\omega} \right), \quad a^{\dagger} = \sqrt{\frac{m\omega}{2\hbar}} \left( x - \frac{ip}{m\omega} \right) \]
    Work out the action of the operators $a$ and $a^{\dagger}$ on the energy eigenstates of the Hamiltonian. \pts{6}
\end{parts}

\medskip\rule{\linewidth}{0.2pt}




\noindent   \textbf{Question 5.}

\begin{parts}
  \item The function $f( \theta,\phi)$ is an eigenfunction of the operator $L_z$ with eigenvalue $m\hbar$, such that $L_+ f( \theta,\phi) = 0$. Find the eigenvalue of the operator $L^2$ acting on the function $f( \theta,\phi)$. \pts{6}
  \item The radial part of the wave function in the first excited state of the Hydrogen atom ($2p$, $l=1$, $m=0$) is given by:
    \[ R_{21}(r) = C r e^{-r/2a} \]
    where $a$ is the Bohr radius. Find the normalization constant $C$ and the most probable distance of the electron from the origin. \pts{8}
\end{parts}

\vfill
\rule{\linewidth}{0.4pt}

\begin{center}\small
  \href{https://bhu-exam.vercel.app/}{Click here to visit our website}\\[4pt]
  \href{https://me-aryan.vercel.app/}{Click here to visit our contributor}\\[4pt]
  \href{https://quashan-me.vercel.app/}{Click here to visit our contributor}
\end{center}

\end{document}
